\section{Experiments and failure attribution}
\label{sec:C09:experiments}
Define comparison axes before allocating resources: single-attempt success, coverage at a candidate budget, selected-answer accuracy, repeated-run consistency, verified constraint violations, latency tails, and complete cost. Include a direct-answer baseline, ordinary intermediate reasoning, independent sampling, aggregation, and search with a common verifier where feasible. A stronger verifier is a meaningful system improvement, but should be named as the intervention. Oracle selection can estimate generation headroom and must be labeled nondeployable. It should not appear in the main autonomous-performance ranking.



Failure attribution should retain both candidate-level and task-level evidence. If no correct candidate exists, selection cannot solve the task. If a correct candidate exists but a weak judge rejects it, more generation may be wasteful. If an accepted program passes visible tests but fails a held-out requirement, the issue can lie in specification coverage. If a correct action is generated from a stale context, execution may still fail. These cases require different interventions. A single score labeled reasoning ability cannot identify them.

The practical role of inference-time computation is to convert available model competence and evidence into a better decision under an explicit resource envelope. Its limits arise from shared misconceptions, imperfect selection, stale state, and external effects as much as from insufficient token generation. The most informative studies expose those mechanisms with mathematical assumptions, implementation paths, and controlled outcomes. They establish where additional computation helps, where it merely repeats an error, and where another real observation is more valuable than another internal continuation.

\paragraph{An auditable search trace.}
The checked breadth-first search implementation offers a useful place to instrument these distinctions. Record every proposal before scoring, every value returned by the evaluator, the retained frontier, and the final selection. On a task with an exact checker, label all generated candidates after the run without exposing those labels to search. This yields separate estimates of candidate coverage, evaluator discrimination, pruning loss, and final selection loss. A correct candidate pruned early is different evidence from a candidate that never appears. The instrumentation should include cache hits and repeated proposals so that duplicated work does not inflate apparent search diversity.

For a fixed set of candidates, compare the deployed selector with random choice, majority answer, shortest valid artifact, and oracle choice. These are selector ablations, not alternative generators. Then hold the selector fixed while varying proposal count and temperature. This two-stage design identifies whether a reported gain comes from finding more correct artifacts or recognizing them better. When the task requires external execution, candidate checks should occur in isolated environments or on immutable artifacts. Otherwise the checker itself can alter the state that later candidates encounter, invalidating the fixed-candidate interpretation.

The score-error bound earlier in this chapter also suggests a calibration diagnostic. Partition candidates by whether they were sampled randomly, selected for high score, or generated adaptively after feedback. Estimate judge error separately in each stratum. A much larger error in the selected stratum is direct evidence that ordinary calibration does not survive optimization pressure. Recalibrating only on the selected examples can help within that regime, but a new adaptive optimizer can create another shift. Keep an untouched evaluation sample and record each judge revision rather than treating the judge as a stationary oracle.

Memory reuse should be audited at the assertion level. If three workers all repeat a conclusion derived from one unverified artifact, their agreement is one evidence lineage, not three independent observations. A provenance graph can collapse these dependent claims before aggregation. Conversely, two independent executable checks on different relevant cases provide distinct evidence even when their natural-language summaries are identical. The controller should reason about that dependence explicitly. Otherwise increasing the number of agents can make a weak claim appear more certain through repetition alone.

These controls support practical allocation decisions. If generation coverage saturates but selection remains weak, spend budget on a stronger checker or a more informative test. If the checker is reliable but no correct candidate appears, improve decomposition, evidence retrieval, or the model. If both are adequate but state changes invalidate the chosen action, reduce observation age and strengthen execution preconditions. The resulting diagnosis connects inference-time reasoning to the surrounding agent system instead of using additional tokens as the default response to every failure.
